\documentclass[a4paper,11pt]{article}

% ---------- Packages ----------
\usepackage[ a4paper, top=0.55in, bottom=0.55in, left=0.65in, right=0.65in ]{geometry}
\usepackage{titlesec}
\usepackage{enumitem}
\usepackage[hidelinks]{hyperref}
\usepackage{xcolor}
\usepackage{tabularx}
\usepackage{fontawesome5}
\usepackage{array}
\usepackage{multicol}
\usepackage{ragged2e}
\usepackage{graphicx}
\usepackage{wrapfig}

% ---------- Font ----------
\renewcommand{\familydefault}{\sfdefault}

% ---------- Colors ----------
\definecolor{primary}{RGB}{30,30,30}
\definecolor{secondary}{RGB}{80,80,80}

% ---------- Hyperlink Setup ----------
\hypersetup{
    colorlinks=true,
    urlcolor=black
}

% ---------- Remove Page Numbers ----------
\pagestyle{empty}

% ---------- Section Formatting ----------
\titleformat{\section}
{\large\bfseries\color{primary}}
{}
{0em}
{}[\titlerule]

\titlespacing{\section}{0pt}{5pt}{4pt}

% ---------- List Spacing ----------
\setlist[itemize]{
    noitemsep,
    topsep=1pt,
    parsep=0pt,
    partopsep=0pt,
    leftmargin=16pt
}

% ---------- Custom Commands ----------
\newcommand{\resumeSubheading}[4]{
\textbf{#1} \hfill \textit{#2} \\
\textit{#3} \hfill \textit{#4}
\vspace{4pt}
}

\newcommand{\resumeProject}[2]{
\textbf{#1} \hfill {\small #2}
}

% ---------- Document ----------
\begin{document}



% ---------- Header ----------
% ---------- ELITE HEADER UI V2 ----------

\begin{center}

{\Huge \textbf{Rishvin Reddy}}\\[4pt]

{\normalsize
Computer Science Engineering Student
\textbar\
Cybersecurity
\textbar\
IoT
\textbar\
Software Systems
}

\vspace{6pt}

\small

\faPhone\ +91 98487 23235
\hspace{1.2em}
\faEnvelope\
\href{mailto:rishvinreddy@gmail.com}
{rishvinreddy@gmail.com}
\hspace{1.2em}
\faMapMarker*\ Hyderabad, India

\vspace{4pt}

\href{https://github.com/RishvinReddy}
{\textbf{GitHub}}
\hspace{1.5em}
\href{https://rishvinreddy.github.io}
{\textbf{Portfolio}}
\hspace{1.5em}
\href{https://www.linkedin.com/in/rishvinreddy/}
{\textbf{LinkedIn}}

\vspace{6pt}
\hrule

\end{center}

\vspace{-8pt}


% ---------- Professional Summary ----------

\section{Professional Summary}

Computer Science Engineering student specializing in \textbf{Blockchain, IoT, and Cybersecurity (BIC)} with strong interests in \textbf{Cybersecurity Engineering, IoT Systems, Embedded Automation, Secure Architectures, and Software Development}. Hands-on experience in developing intelligent automation systems, embedded IoT applications, secure authentication workflows, and software-driven engineering solutions through academic, personal, and project-based work.

Possesses practical experience in building \textbf{IoT-enabled smart systems, biometric authentication platforms, operating systems simulations, full-stack applications, and algorithmic problem-solving systems}. Strong technical foundation in \textbf{Python, JavaScript, Data Structures \& Algorithms, Operating Systems, Database Management Systems, Web Technologies, and Embedded System Integration} using platforms such as Arduino, ESP32, Flask, React, and PostgreSQL.

Actively involved in engineering projects focused on \textbf{smart automation, secure device-to-cloud communication, intelligent monitoring systems, biometric security, systems simulation, and real-world problem solving}. Demonstrates a growing interest in \textbf{secure software engineering, intelligent connected systems, scalable architectures, cybersecurity practices, and AI-assisted engineering workflows} for solving practical and large-scale technical challenges.

Strongly motivated to build \textbf{secure, scalable, and intelligent systems} combining cybersecurity, IoT, software engineering, and intelligent automation with a long-term vision toward engineering impactful real-world technologies and scalable products.


% ---------- Education ----------

\section{Education}

\resumeSubheading
{Woxsen University}{Hyderabad, India}
{Bachelor of Technology (B.Tech) in Computer Science and Engineering\\
Specialization: Blockchain, IoT \& Cybersecurity (BIC)}
{2024 -- 2028}

\begin{itemize}
    \item \textbf{CGPA:} 9.01/10
    \item Relevant Coursework: Data Structures \& Algorithms, Object-Oriented Programming, Database Management Systems, Web Technologies, IoT Engineering, Computational Thinking, Probability \& Statistics
\end{itemize}


% ---------- Technical Skills ----------

\section{Technical Skills}

\begin{itemize}[leftmargin=0.15in]

\item \textbf{Programming Languages:}
Python, JavaScript, SQL, HTML, CSS

\item \textbf{Frameworks \& Libraries:}
React, Flask

\item \textbf{Web Technologies:}
Frontend Development, Responsive Web Design, UI Development, Form Handling, Dashboard Interfaces

\item \textbf{IoT \& Embedded Systems:}
Arduino UNO, ESP32, Tinkercad, Sensor Integration, Embedded Systems, Serial Communication, Smart Automation, Real-Time Monitoring, Device Interfacing, Telemetry Systems

\item \textbf{Cybersecurity \& Secure Systems:}
Cybersecurity Fundamentals, Network Security, Authentication Systems, Secure Development Practices, Threat Analysis Basics, Biometric Authentication, Identity Verification, Secure Device-to-Cloud Communication, Security-Oriented System Design

\item \textbf{Databases \& Storage:}
SQL, PostgreSQL, Database Management Systems (DBMS), Data Persistence, Structured Data Management

\item \textbf{Developer Tools \& Platforms:}
Git, GitHub, Antigravity IDE, GitHub Pages, Tinkercad, Arduino Ecosystem, Overleaf, Coursera Learning Platform

\item \textbf{Software Engineering Concepts:}
Object-Oriented Programming (OOP), Data Structures \& Algorithms (DSA), Problem Solving, Algorithm Design, Debugging, Software Development Fundamentals, Application Development

\item \textbf{Systems Engineering Concepts:}
Operating Systems, Memory Management, Disk Scheduling Algorithms, Systems Simulation, Performance Analysis, Computational Thinking

\item \textbf{Engineering Methodologies \& Practices:}
Project-Based Learning, Prototype Development, Engineering Problem Solving, Technical Documentation, System Design Thinking, Workflow Optimization

\end{itemize}




% ---------- Domain Expertise ----------

\section{Domain Expertise}

\begin{itemize}[leftmargin=0.15in]

\item \textbf{Cybersecurity Engineering}
\begin{itemize}
    \item Authentication Systems
    \item Biometric Verification
    \item Secure Development Practices
    \item Identity Validation
    \item Threat Analysis Fundamentals
    \item Network Security Basics
    \item Secure Device-to-Cloud Communication
    \item Security-Oriented System Design
\end{itemize}

\item \textbf{Internet of Things (IoT) \& Embedded Systems}
\begin{itemize}
    \item Arduino UNO \& ESP32 Development
    \item Sensor Integration \& Hardware Interfacing
    \item Embedded Automation Systems
    \item Smart Agriculture Systems
    \item Smart Monitoring Applications
    \item Telemetry \& Real-Time Monitoring
    \item Device Communication Workflows
    \item Prototype-Based IoT Development
\end{itemize}

\item \textbf{Software Engineering \& Application Development}
\begin{itemize}
    \item Full-Stack Development Fundamentals
    \item Frontend Development
    \item Dashboard-Based Interfaces
    \item Database Integration
    \item Application Workflow Design
    \item Software Problem Solving
    \item Debugging \& Optimization
    \item Scalable System Thinking
\end{itemize}

\item \textbf{Systems Engineering \& Algorithmic Computing}
\begin{itemize}
    \item Data Structures \& Algorithms
    \item Pattern Matching Algorithms
    \item Operating Systems Concepts
    \item Memory Management Simulation
    \item Disk Scheduling Algorithms
    \item Performance Analysis
    \item Computational Thinking
    \item Systems-Level Problem Solving
\end{itemize}

\item \textbf{Engineering Design \& Innovation}
\begin{itemize}
    \item Patent-Oriented Thinking
    \item Intelligent System Design
    \item Prototype Development
    \item Secure Architecture Thinking
    \item Automation-Driven Engineering
    \item Product-Oriented Development
    \item Technical Documentation
    \item Engineering Workflow Optimization
\end{itemize}

\end{itemize}






% ---------- Featured Engineering Projects ----------

\section{Featured Engineering Projects}

\resumeProject
{IoT-Enabled Smart Irrigation System for Soil Health Monitoring}
{Python, Arduino UNO, ESP32, Embedded Systems, IoT, Sensor Automation}

\textbf{Project Overview:}
Developed an intelligent IoT-enabled smart irrigation system focused on soil health monitoring, automated irrigation decision-making, and real-time environmental monitoring for smart agriculture applications.

\textbf{Problem Statement:}
Traditional irrigation methods often lead to water wastage, inefficient monitoring, and lack of automation. This project addresses the need for intelligent, sensor-driven irrigation systems capable of optimizing water utilization and improving agricultural efficiency.

\textbf{System Architecture:}

Soil Moisture Sensor / Temperature Sensor
$\rightarrow$ Arduino UNO Processing Layer
$\rightarrow$ Threshold-Based Decision Engine
$\rightarrow$ Servo Motor Irrigation Control
$\rightarrow$ ESP32 IoT Layer
$\rightarrow$ Real-Time Dashboard Monitoring

\textbf{Tech Stack:}
Python, Arduino UNO, ESP32, Embedded C/C++, Tinkercad, Sensor Integration, Serial Communication, IoT Dashboard

\textbf{Core Features:}
\begin{itemize}
    \item Real-time soil moisture monitoring using sensor-driven environmental data collection.
    
    \item Automated irrigation control using threshold-based decision logic and servo motor automation.
    
    \item Multi-level LED indicators to represent soil moisture conditions ranging from dry to optimal levels.
    
    \item ESP32-enabled IoT dashboard integration for remote environmental monitoring and telemetry visualization.
    
    \item Smart agriculture workflow designed to reduce water wastage and improve irrigation efficiency.
\end{itemize}

\textbf{Engineering Concepts Applied:}
Embedded Systems, IoT Automation, Smart Agriculture, Sensor Integration, Real-Time Monitoring, Decision Logic Systems, Telemetry, Device Communication.

\textbf{Role \& Contribution:}
Designed, developed, and implemented the complete automation workflow, sensor integration logic, irrigation decision system, and IoT monitoring pipeline.

\textbf{Future Scope:}
Potential integration with predictive analytics, weather forecasting APIs, cloud-based irrigation recommendations, and mobile application monitoring systems.



\resumeProject
{VoteSafe -- Secure Biometric Voting System}
{React 19, Vite, Node.js, Express.js, PostgreSQL, AES-256-GCM, Biometrics}

\textbf{Project Overview:}
Designed and developed a secure biometric voting platform focused on authenticated digital voting, voter identity validation, secure data handling, and tamper-resistant voting workflows using modern full-stack technologies and biometric hardware integration.

\textbf{Problem Statement:}
Traditional and digital voting systems face challenges including identity fraud, duplicate voting, insecure authentication, lack of transparency, and centralized trust limitations. VoteSafe was developed to improve election security through biometric authentication, encrypted identity verification, and auditable voting workflows.

\textbf{System Architecture:}

Fingerprint Capture (Mantra MFS110 Scanner)
$\rightarrow$ RD Service Polling
$\rightarrow$ ISO Template Extraction
$\rightarrow$ AES-256-GCM Encryption
$\rightarrow$ PostgreSQL Secure Match Validation
$\rightarrow$ Vote Authorization
$\rightarrow$ Merkle Audit Ledger Logging

\textbf{Tech Stack:}
React 19, Vite, Node.js, Express.js, PostgreSQL, AES-256-GCM Encryption, Mantra MFS110 Fingerprint Scanner, RD Service Polling, Secure Template Validation

\textbf{Core Features:}
\begin{itemize}

    \item Fingerprint-based biometric authentication using Mantra MFS110 hardware for secure voter verification.

    \item Secure ISO biometric template extraction and encrypted identity matching workflow.

    \item AES-256-GCM encryption for protecting biometric information and secure voter authentication pipelines.

    \item Duplicate vote prevention mechanism through secure identity validation and authorization logic.

    \item Merkle-tree based audit ledger for tamper-resistant vote activity tracking and verification.

    \item Session-isolated public voter mode for improved voting privacy and system security.

    \item Full-stack application architecture supporting secure frontend-backend communication and database integration.

\end{itemize}

\textbf{Cybersecurity Concepts Applied:}
Biometric Authentication, Secure Identity Verification, Encryption Standards, Authentication Workflows, Secure Session Handling, Data Protection, Access Validation, Security-Oriented System Design.

\textbf{Software Engineering Concepts Applied:}
Frontend-Backend Integration, Database Systems, RESTful Workflow Architecture, State Management, System Design, Secure API Communication.

\textbf{Role \& Contribution:}
Independently designed and developed the complete system architecture, frontend-backend workflow, biometric verification pipeline, encryption workflow, secure authentication process, and database integration.

\textbf{Future Scope:}
Potential expansion into blockchain-based vote immutability, distributed verification mechanisms, cloud-scale deployment, advanced identity systems, and national-scale secure election infrastructure.




\resumeProject
{IoT Smart Waste Monitoring System}
{Arduino UNO, Embedded Systems, Ultrasonic Sensor, Serial Communication, Smart Monitoring}

\textbf{Project Overview:}
Developed an IoT-based smart waste monitoring system focused on intelligent garbage level detection, automated status monitoring, and efficient waste management workflows through embedded sensing and device communication.

\textbf{Problem Statement:}
Traditional waste collection systems often suffer from inefficient scheduling, unnecessary collection cycles, and lack of visibility into bin fill status. This project addresses these inefficiencies through sensor-driven waste monitoring and intelligent fill-level detection.

\textbf{System Architecture:}

Ultrasonic Sensor
$\rightarrow$ Sender Arduino UNO
$\rightarrow$ Fill-Level Threshold Logic
$\rightarrow$ Serial Communication Pipeline
$\rightarrow$ Receiver Arduino UNO
$\rightarrow$ LED Status Display / Monitoring Output

\textbf{Tech Stack:}
Arduino UNO, Embedded C/C++, Ultrasonic Sensor, Tinkercad, Serial Communication, Embedded Systems

\textbf{Core Features:}
\begin{itemize}

    \item Real-time garbage fill-level detection using ultrasonic sensing for intelligent waste monitoring.

    \item Threshold-based fullness evaluation to identify different bin capacity levels and trigger system responses.

    \item LED-based visual status indication for monitoring waste levels and operational visibility.

    \item Arduino-to-Arduino serial communication enabling remote status transmission and distributed monitoring.

    \item Smart waste management workflow designed to improve collection efficiency and reduce unnecessary servicing cycles.

\end{itemize}

\textbf{Engineering Concepts Applied:}
IoT Automation, Embedded Systems, Smart Monitoring, Serial Communication, Sensor Integration, Real-Time Data Processing, Device Communication, Automation Logic.

\textbf{Role \& Contribution:}
Collaborated in the design and development of the embedded monitoring workflow, sensor threshold logic, communication architecture, and intelligent monitoring system for waste management applications.

\textbf{Future Scope:}
Potential integration with cloud dashboards, GPS-enabled waste collection optimization, predictive analytics for collection planning, and smart city infrastructure systems.


\resumeProject
{Text Search Engine Using Trie and Pattern Matching Algorithms}
{Python, Data Structures \& Algorithms, Pattern Matching, Search Optimization}

\textbf{Project Overview:}
Developed an efficient text search engine focused on optimized keyword retrieval, fast pattern matching, and comparative search performance analysis using advanced data structures and algorithmic techniques.

\textbf{Problem Statement:}
Traditional brute-force text searching becomes inefficient for large datasets due to higher computational overhead and slower search performance. This project was developed to improve text retrieval efficiency through optimized indexing and pattern matching strategies.

\textbf{System Architecture:}

Input Text / Document Data
$\rightarrow$ Trie-Based Indexing
$\rightarrow$ Pattern Matching Engine
(KMP + Brute Force)
$\rightarrow$ Search Processing
$\rightarrow$ Optimized Result Retrieval

\textbf{Tech Stack:}
Python, Trie Data Structure, KMP (Knuth--Morris--Pratt), Brute Force String Matching, Algorithm Analysis

\textbf{Core Features:}
\begin{itemize}

    \item Implemented Trie data structure for efficient keyword indexing and optimized text retrieval.

    \item Developed KMP (Knuth--Morris--Pratt) pattern matching algorithm for improved substring search performance.

    \item Implemented brute-force string matching to compare traditional search performance against optimized approaches.

    \item Performed comparative algorithm analysis to evaluate time complexity, search efficiency, and retrieval performance.

    \item Designed an optimized search workflow for faster keyword lookup and intelligent text processing.

\end{itemize}

\textbf{Engineering Concepts Applied:}
Data Structures \& Algorithms (DSA), Pattern Matching, Search Optimization, Time Complexity Analysis, Computational Thinking, Efficient Data Retrieval.

\textbf{Software Engineering Concepts Applied:}
Algorithm Design, Backend Logic Development, Performance Optimization, Search Workflow Design, Comparative Analysis.

\textbf{Role \& Contribution:}
Contributed as Backend Developer in Team 11 by implementing search logic, algorithm workflows, Trie-based indexing, and performance-oriented retrieval mechanisms.

\textbf{Future Scope:}
Potential integration with NLP pipelines, semantic search, large-scale indexing systems, intelligent document retrieval, and AI-assisted search optimization systems.






\resumeProject
{Interactive Disk Scheduling Algorithm Visualizer}
{Python, Operating Systems, Algorithm Simulation, Performance Analysis}

\textbf{Project Overview:}
Developed an interactive operating systems simulation platform for visualizing and analyzing disk scheduling algorithms through real-time head movement simulation, seek time computation, and comparative performance evaluation.

\textbf{Problem Statement:}
Disk scheduling plays a critical role in operating system efficiency and storage performance. Traditional theoretical learning often lacks practical visualization of algorithm behavior, seek optimization, and performance trade-offs. This project addresses these limitations through an interactive simulation environment.

\textbf{System Architecture:}

User Input Queue
$\rightarrow$ Scheduling Algorithm Selection
(FCFS / SSTF / SCAN / C-SCAN / LOOK / C-LOOK)
$\rightarrow$ Disk Head Movement Simulation
$\rightarrow$ Seek Time Calculation
$\rightarrow$ Comparative Visualization Output

\textbf{Tech Stack:}
Python, Operating Systems Concepts, Algorithm Simulation, Scheduling Logic, Performance Analysis

\textbf{Algorithms Implemented:}
\begin{itemize}
    \item FCFS (First Come First Serve)
    \item SSTF (Shortest Seek Time First)
    \item SCAN
    \item C-SCAN
    \item LOOK
    \item C-LOOK
\end{itemize}

\textbf{Core Features:}
\begin{itemize}

    \item Implemented multiple disk scheduling algorithms for comparative operating systems performance evaluation.

    \item Developed disk head movement simulation to visualize request servicing behavior and optimization strategies.

    \item Calculated seek time performance for evaluating efficiency across different scheduling approaches.

    \item Designed an interactive learning environment for understanding disk scheduling behavior and algorithm trade-offs.

    \item Performed comparative analysis to study efficiency, optimization, and scheduling performance under varying request sequences.

\end{itemize}

\textbf{Engineering Concepts Applied:}
Operating Systems, Scheduling Algorithms, Performance Optimization, Systems Simulation, Computational Thinking, Comparative Analysis.

\textbf{Software Engineering Concepts Applied:}
Algorithm Design, Interactive Simulation Development, Visualization Workflows, Logic Implementation, Performance Evaluation.

\textbf{Role \& Contribution:}
Designed and implemented algorithm workflows, scheduling simulation logic, comparative evaluation mechanisms, and performance analysis features.

\textbf{Future Scope:}
Potential integration with GUI-based visualization systems, advanced storage scheduling analysis, performance benchmarking, and educational operating systems simulators.




\resumeProject
{OS Paging and Virtual Memory Simulation}
{Python, Operating Systems, Memory Management, Algorithm Simulation}

\textbf{Project Overview:}
Developed an operating systems simulation platform focused on virtual memory management, page replacement optimization, and comparative analysis of memory allocation strategies through interactive paging workflows.

\textbf{Problem Statement:}
Efficient memory management is a critical challenge in operating systems due to limited physical memory and varying process demands. Traditional learning methods often fail to provide practical understanding of page replacement behavior and performance trade-offs. This project addresses these limitations through an interactive simulation of virtual memory mechanisms.

\textbf{System Architecture:}

Page Reference String
$\rightarrow$ Memory Frame Allocation
$\rightarrow$ Page Replacement Algorithm Selection
(FIFO / LRU / Optimal)
$\rightarrow$ Page Fault Tracking
$\rightarrow$ Comparative Performance Evaluation

\textbf{Tech Stack:}
Python, Operating Systems Concepts, Virtual Memory Management, Page Replacement Algorithms, Systems Simulation

\textbf{Algorithms Implemented:}
\begin{itemize}
    \item FIFO (First In First Out)
    \item LRU (Least Recently Used)
    \item Optimal Page Replacement
\end{itemize}

\textbf{Core Features:}
\begin{itemize}

    \item Implemented FIFO, LRU, and Optimal page replacement algorithms for comparative virtual memory analysis.

    \item Simulated memory frame allocation and replacement workflows to study page replacement behavior.

    \item Developed page fault tracking mechanisms to evaluate algorithm efficiency under varying memory access patterns.

    \item Designed an interactive simulation environment for understanding virtual memory concepts and replacement strategies.

    \item Performed comparative performance analysis to identify efficiency trade-offs among paging algorithms.

\end{itemize}

\textbf{Engineering Concepts Applied:}
Operating Systems, Virtual Memory Management, Paging Mechanisms, Algorithm Simulation, Performance Analysis, Systems-Level Problem Solving.

\textbf{Software Engineering Concepts Applied:}
Simulation Workflow Design, Algorithm Implementation, Logic Development, Comparative Evaluation, Computational Thinking.

\textbf{Role \& Contribution:}
Designed and implemented paging logic, algorithm workflows, page replacement simulation, and comparative memory management analysis mechanisms.

\textbf{Future Scope:}
Potential integration with graphical visualization systems, advanced operating systems simulation environments, educational learning platforms, and memory performance benchmarking systems.









% ---------- Patent & Innovation ----------

\section{Patent \& Innovation}

\resumeProject
{Government of India Registered Design Patent -- IoT Connectivity Device}
{Design No: 470097-001 \textbar\ Registered: Aug 2025 \textbar\ Status: Active}

\textbf{Innovation Overview:}
Co-Inventor of a Government of India registered hardware design patent focused on secure, authenticated, and intelligent device-to-cloud communication for IoT-enabled systems.

\textbf{Problem Statement:}
Traditional IoT connectivity systems often face challenges related to secure authentication, telemetry handling, reliable device communication, and scalable real-time monitoring. This innovation addresses these challenges through an intelligent IoT hardware interface design.

\textbf{Patent Scope \& Technical Focus:}

Intelligent IoT Device Interface
$\rightarrow$ Secure Device Authentication
$\rightarrow$ Real-Time Telemetry Transmission
$\rightarrow$ Device-to-Cloud Communication
$\rightarrow$ Monitoring \& Data Exchange

\textbf{Technical Significance:}
\begin{itemize}

    \item Designed to enable secure and authenticated IoT device communication workflows.

    \item Focused on improving telemetry transmission and intelligent connected-system monitoring.

    \item Supports real-time communication between devices and cloud-connected infrastructure.

    \item Demonstrates practical innovation in embedded systems and secure IoT architectures.

\end{itemize}

\textbf{Engineering Concepts Applied:}
IoT Systems, Secure Device Communication, Telemetry Systems, Embedded Architecture Thinking, Real-Time Monitoring, Intelligent Connected Systems.

\textbf{Role \& Contribution:}
Contributed as \textbf{Co-Inventor} in the ideation, design thinking, and conceptualization of the patented IoT connectivity hardware interface and secure communication workflow.

\textbf{Patent Information:}
\begin{itemize}
    \item \textbf{Design Patent Office:} Government of India
    
    \item \textbf{Patent Category:} IoT System
    
    \item \textbf{Patent Status:} Registered \& Active
    
    \item \textbf{Design Number:} 470097-001
\end{itemize}

\textbf{Future Scope:}
Potential applications in smart monitoring systems, industrial IoT, intelligent automation systems, secure telemetry infrastructure, and scalable connected-device ecosystems.








% ---------- Certifications ----------

\section{Certifications}

\textbf{Cybersecurity \& Secure Systems}
\begin{itemize}
    \item \textbf{Introduction to Cyber Attacks} --- Coursera / New York University (2024)
    
    \begin{itemize}
        \item Covered cybersecurity fundamentals including threat modeling, attack vectors, vulnerability analysis, and secure systems concepts.
    \end{itemize}
\end{itemize}

\vspace{3pt}

\textbf{Programming \& Software Engineering}
\begin{itemize}

    \item \textbf{Programming for Everybody (Getting Started with Python)} --- Coursera / University of Michigan (2023)

    \begin{itemize}
        \item Learned foundational Python programming concepts including variables, functions, loops, expressions, and logical problem solving.
    \end{itemize}

    \item \textbf{Data Structures Using Python --- An Introduction} --- Coursera (2023)

    \begin{itemize}
        \item Explored core data structures including lists, stacks, queues, trees, and graphs through Python-based implementations.
    \end{itemize}

    \item \textbf{Java Programming} --- Coursera (2023)

    \begin{itemize}
        \item Covered object-oriented programming concepts including classes, inheritance, interfaces, and collections in Java.
    \end{itemize}

    \item \textbf{Object-Oriented Design} --- Coursera / University of Alberta (2023)

    \begin{itemize}
        \item Learned software design principles including SOLID concepts, UML diagrams, design patterns, and software architecture fundamentals.
    \end{itemize}

\end{itemize}

\vspace{3pt}

\textbf{Embedded Systems, IoT \& Hardware Systems}
\begin{itemize}

    \item \textbf{Interfacing with the Arduino} --- Coursera / University of California, Irvine (2023)

    \begin{itemize}
        \item Covered sensor interfacing, actuator integration, and hardware communication using Arduino microcontrollers.
    \end{itemize}

    \item \textbf{The Arduino Platform and C Programming} --- Coursera / University of California, Irvine (2023)

    \begin{itemize}
        \item Learned Arduino platform fundamentals and embedded programming concepts using C-based firmware logic.
    \end{itemize}

    \item \textbf{Nanotechnology and Nanosensors, Part 1} --- Coursera / Technion (2024)

    \begin{itemize}
        \item Studied nanosensor design principles, nanomaterials, and emerging sensing technologies for intelligent systems.
    \end{itemize}

\end{itemize}

\vspace{3pt}

\textbf{Emerging Technologies}
\begin{itemize}

    \item \textbf{Quantum Computing For Everyone --- An Introduction} --- Coursera (2024)

    \begin{itemize}
        \item Explored foundational quantum computing concepts including qubits, superposition, entanglement, and quantum algorithms.
    \end{itemize}

\end{itemize}

\vspace{3pt}

\textbf{Professional Development \& Communication}
\begin{itemize}

    \item \textbf{Write Professional Emails in English} --- Coursera / Georgia Tech (2023)

    \begin{itemize}
        \item Learned principles of professional written communication and structured business correspondence.
    \end{itemize}

    \item \textbf{Understanding and Optimizing Campus Building Energy Consumption} --- Woxsen University (2025)

    \begin{itemize}
        \item Participated in a hands-on technical workshop focused on practical approaches to energy optimization and sustainability systems.
    \end{itemize}

\end{itemize}









% ---------- Additional Engineering Projects ----------

\section{Additional Engineering Projects}

\resumeProject
{Smart Budget Planner}
{TypeScript, Frontend Development, Dashboard Systems}

\textbf{Project Overview:}
Developed an interactive budget planning application focused on expense tracking, spending visualization, financial planning, and budget optimization through an intuitive dashboard interface.

\textbf{Core Features:}
\begin{itemize}
    \item Income and expense tracking workflow for personal financial management.
    
    \item Spending analysis and visualization for monitoring financial behavior.
    
    \item Goal-oriented budget planning for improved financial organization.
    
    \item Interactive dashboard interface for intuitive user experience and financial insights.
\end{itemize}

\textbf{Engineering Concepts Applied:}
Frontend Development, Dashboard Systems, User Interface Design, Data Visualization, Workflow Design.

\vspace{4pt}

\resumeProject
{Government Payroll Management System}
{Python, SQL, Software Systems}

\textbf{Project Overview:}
Developed a payroll automation system focused on employee record management, salary computation, payroll generation, and structured organizational data handling.

\textbf{Core Features:}
\begin{itemize}
    \item Employee record management workflow for structured personnel tracking.
    
    \item Automated salary calculation and payroll processing mechanisms.
    
    \item Payroll reporting and structured data organization for operational management.
    
    \item Database-driven management for payroll-related information processing.
\end{itemize}

\textbf{Engineering Concepts Applied:}
Database Management, Automation Logic, CRUD Operations, Workflow Design, Structured Data Management.

\vspace{4pt}

\resumeProject
{Outing Form Management System}
{HTML, CSS, JavaScript, Workflow Systems}

\textbf{Project Overview:}
Designed a web-based outing management system enabling student request submission, approval workflows, and structured administrative management.

\textbf{Core Features:}
\begin{itemize}
    \item Student outing request submission and structured form validation.
    
    \item Administrative workflow for request monitoring and record management.
    
    \item Dynamic student management for scalable request handling.
    
    \item Structured approval process supporting organized workflow execution.
\end{itemize}

\textbf{Engineering Concepts Applied:}
Web Development, Workflow Automation, Form Validation, User Interface Design, Record Management.

\vspace{4pt}

\resumeProject
{Face Mesh Verification System}
{Python, Computer Vision, Biometric Verification}

\textbf{Project Overview:}
Developed a computer vision-based facial verification system focused on face landmark detection, mesh tracking, and biometric identity validation.

\textbf{Core Features:}
\begin{itemize}
    \item Real-time facial landmark tracking and mesh extraction workflow.
    
    \item Biometric identity verification through facial structure analysis.
    
    \item Computer vision pipeline for image processing and verification workflows.
    
    \item Real-time face analysis for authentication-oriented experimentation.
\end{itemize}

\textbf{Engineering Concepts Applied:}
Computer Vision, Biometric Verification, Face Recognition Concepts, Image Processing, Real-Time Detection Systems.








% ---------- Technical Experience & Leadership ----------

\section{Technical Experience \& Leadership}

\textbf{Backend Developer --- Team 11, Data Structures \& Algorithms Project}
\hfill \textit{Academic Engineering Team Project}

\begin{itemize}

    \item Contributed as Backend Developer in the development of a \textbf{Text Search Engine using Trie and Pattern Matching Algorithms}.

    \item Implemented backend logic involving \textbf{Trie-based indexing, KMP pattern matching, brute-force search mechanisms, and performance-oriented search optimization}.

    \item Participated in algorithm implementation, search workflow design, comparative performance analysis, and computational optimization.

    \item Collaborated with team members to improve system efficiency, retrieval logic, and technical implementation workflows.

\end{itemize}

\vspace{6pt}

\textbf{Co-Inventor --- Government of India Registered Design Patent}
\hfill \textit{IoT System Innovation}

\begin{itemize}

    \item Contributed as \textbf{Co-Inventor} in the ideation, design thinking, and technical conceptualization of a Government of India registered IoT hardware design patent.

    \item Participated in the development of an intelligent IoT interface focused on \textbf{secure device-to-cloud communication, telemetry workflows, and connected-system architecture}.

    \item Worked on innovation-focused problem solving involving \textbf{IoT connectivity, secure communication logic, and real-time system interaction}.

\end{itemize}

\vspace{6pt}

\textbf{Independent Technical Project Development}
\hfill \textit{Personal \& Academic Engineering Work}

\begin{itemize}

    \item Independently designed and developed multiple engineering projects spanning \textbf{Cybersecurity, IoT, Embedded Systems, Software Engineering, Computer Vision, and Systems Simulation}.

    \item Built practical systems involving \textbf{biometric authentication, IoT automation, operating systems simulation, smart monitoring, and secure software workflows}.

    \item Applied engineering concepts including \textbf{system design, algorithm optimization, secure architecture thinking, automation logic, workflow management, and performance analysis}.

    \item Demonstrated strong initiative in learning, building, experimentation, and implementation of real-world engineering solutions.

\end{itemize}

\vspace{6pt}

\textbf{Technical Collaboration \& Engineering Participation}

\begin{itemize}

    \item Collaborated on project-based engineering activities involving \textbf{IoT systems, embedded automation, algorithmic problem solving, and software development workflows}.

    \item Participated in hands-on academic and technical environments emphasizing \textbf{problem solving, prototyping, experimentation, and practical implementation}.

\end{itemize}






% ---------- Research & Engineering Interests ----------

\section{Research \& Engineering Interests}

\begin{itemize}[leftmargin=0.15in]

\item \textbf{Cybersecurity Engineering}
\begin{itemize}
    \item Authentication Systems
    \item Biometric Security
    \item Secure Identity Verification
    \item Security-Oriented System Design
    \item Threat Analysis \& Secure Architectures
\end{itemize}

\item \textbf{Internet of Things (IoT) \& Embedded Systems}
\begin{itemize}
    \item Smart Automation Systems
    \item Intelligent Monitoring Systems
    \item Device-to-Cloud Communication
    \item Sensor-Based Automation
    \item Embedded System Integration
\end{itemize}

\item \textbf{Software Engineering \& Intelligent Systems}
\begin{itemize}
    \item Full-Stack Application Development
    \item Scalable Software Architectures
    \item Intelligent Workflow Systems
    \item Product-Oriented Engineering
    \item Software Problem Solving
\end{itemize}

\item \textbf{Systems Engineering \& Computational Thinking}
\begin{itemize}
    \item Operating Systems Simulation
    \item Performance Optimization
    \item Algorithmic Problem Solving
    \item Memory Management Systems
    \item Efficient Data Processing
\end{itemize}

\item \textbf{Emerging Technologies \& Innovation}
\begin{itemize}
    \item Intelligent Connected Systems
    \item Secure IoT Architectures
    \item AI-Assisted Engineering Workflows
    \item Smart Infrastructure Systems
    \item Scalable Technology Innovation
\end{itemize}

\end{itemize}



% ---------- Platforms & Technical Profiles ----------

\section{Platforms \& Technical Profiles}

\begin{itemize}[leftmargin=0.15in]

\item \textbf{GitHub Profile}
\begin{itemize}
    \item Repository portfolio featuring projects in \textbf{Cybersecurity, IoT, Embedded Systems, Operating Systems, Algorithms, Full-Stack Development, and Intelligent Automation}.
    
    \item Focused on practical engineering implementation, experimentation, project documentation, and software development workflows.
    
    \item Profile: \href{https://github.com/RishvinReddy}{github.com/RishvinReddy}
\end{itemize}

\item \textbf{Portfolio Website}
\begin{itemize}
    \item Personal engineering portfolio showcasing \textbf{technical projects, engineering domains, patent work, software systems, and IoT-focused development}.
    
    \item Demonstrates technical identity, engineering interests, project architecture, and practical system-building experience.
    
    \item Website: \href{https://rishvinreddy.github.io}{rishvinreddy.github.io}
\end{itemize}

\item \textbf{LinkedIn Professional Profile}
\begin{itemize}
    \item Professional networking profile focused on \textbf{Cybersecurity, IoT, Software Engineering, Intelligent Systems, and Engineering Development}.
    
    \item Used for professional networking, technical positioning, and engineering portfolio visibility.
    
    \item Profile: \href{https://www.linkedin.com/in/rishvinreddy/}{linkedin.com/in/rishvinreddy}
\end{itemize}

\item \textbf{Coursera Learning Profile}
\begin{itemize}
    \item Active learning platform used for continuous development in \textbf{Cybersecurity, Software Engineering, Programming, IoT, and Emerging Technologies}.
\end{itemize}

\item \textbf{Engineering Development Ecosystem}
\begin{itemize}
    \item \textbf{Antigravity IDE} for software development and engineering experimentation.
    
    \item \textbf{GitHub Pages} for deployment and portfolio hosting.
    
    \item \textbf{Tinkercad} for IoT simulation, prototyping, and embedded system experimentation.
    
    \item \textbf{Overleaf} for technical writing, engineering documentation, and LaTeX-based professional resume creation.
\end{itemize}

\end{itemize}




% ---------- Relevant Coursework ----------

\section{Relevant Coursework}

\textbf{Computer Science \& Software Engineering}
\begin{itemize}
    \item Data Structures \& Algorithms (DSA)
    \item Object-Oriented Programming (OOP)
    \item Database Management Systems (DBMS)
    \item Operating Systems
    \item Computer Networks
    \item Computational Thinking \& Problem Solving
    \item Web Technologies
    \item Software Development Fundamentals
    \item Technical Writing
\end{itemize}

\vspace{4pt}

\textbf{IoT, Embedded Systems \& Engineering}
\begin{itemize}
    \item Engineering Concepts in IoT
    \item Embedded System Concepts
    \item Sensor-Based System Integration
    \item Smart Agriculture \& Automation Concepts
    \item Real-Time Monitoring Systems
\end{itemize}

\vspace{4pt}

\textbf{Systems Engineering \& Computational Concepts}
\begin{itemize}
    \item Design \& Analysis of Algorithms
    \item Memory Management Concepts
    \item Disk Scheduling Algorithms
    \item Systems Simulation
    \item Performance Analysis
\end{itemize}

\vspace{4pt}

\textbf{Mathematics, Science \& Engineering Foundations}
\begin{itemize}
    \item Engineering Mathematics
    \item Probability \& Statistics
    \item Physics for Engineers
    \item Wave Optics
    \item Quantum Physics
\end{itemize}

\vspace{4pt}

\textbf{Emerging Technologies \& Innovation}
\begin{itemize}
    \item Blockchain Fundamentals
    \item Cybersecurity Fundamentals
    \item Intelligent Systems Concepts
    \item Secure System Thinking
\end{itemize}






% ---------- Workshops & Technical Participation ----------

\section{Workshops \& Technical Participation}

\textbf{Technical Workshops \& Academic Participation}
\begin{itemize}

    \item \textbf{Understanding and Optimizing Campus Building Energy Consumption}
    \hfill \textit{Woxsen University | 2025}

    \begin{itemize}
        \item Participated in a technical workshop focused on energy optimization, sustainability systems, and analytical approaches for reducing energy consumption in institutional environments.
        
        \item Explored practical engineering concepts involving monitoring systems, optimization workflows, and sustainable infrastructure thinking.
    \end{itemize}

\end{itemize}

\vspace{4pt}

\textbf{Engineering Project Participation}
\begin{itemize}

    \item Actively participated in project-based learning involving \textbf{IoT systems, embedded automation, software engineering, operating systems simulation, and cybersecurity-oriented workflows}.

    \item Engaged in hands-on implementation of \textbf{Arduino systems, ESP32 integration, sensor automation, algorithmic problem solving, and intelligent monitoring systems}.

    \item Participated in practical experimentation using \textbf{Tinkercad simulations, embedded hardware integration, and software prototyping workflows}.

\end{itemize}

\vspace{4pt}

\textbf{Technical Learning \& Continuous Development}
\begin{itemize}

    \item Continuously engaged in structured technical learning through \textbf{Coursera, engineering experimentation, project implementation, GitHub-based development, and self-driven technical exploration}.

    \item Regularly explores concepts related to \textbf{Cybersecurity, IoT, Full-Stack Engineering, Systems Engineering, Secure Architectures, and Intelligent Automation Systems}.

\end{itemize}

\vspace{4pt}

\textbf{Innovation \& Engineering Exposure}
\begin{itemize}

    \item Participated in innovation-oriented engineering activities through contribution to a \textbf{Government of India Registered Design Patent} in IoT connectivity systems.

    \item Engaged in applied engineering problem solving through intelligent automation, secure systems, monitoring systems, and embedded technology experimentation.

\end{itemize}




% ---------- Achievements ----------

\section{Achievements}

\begin{itemize}

    \item \textbf{Government of India Registered Design Patent}
    
    Recognized as \textbf{Co-Inventor} of an active Government of India registered IoT hardware design patent (\textbf{Design No: 470097-001}) focused on secure device-to-cloud connectivity and telemetry systems.

    \item \textbf{Strong Academic Performance}
    
    Maintaining a current \textbf{CGPA of 9.01/10} in B.Tech Computer Science and Engineering with specialization in \textbf{Blockchain, IoT \& Cybersecurity (BIC)} at Woxsen University.

    \item \textbf{Multi-Domain Engineering Project Portfolio}
    
    Successfully designed and developed projects spanning \textbf{Cybersecurity, IoT, Embedded Systems, Software Engineering, Operating Systems, Computer Vision, and Intelligent Automation}.

    \item \textbf{Backend Development Contribution}
    
    Contributed as \textbf{Backend Developer} in Team 11 for the implementation of a \textbf{Text Search Engine using Trie and Pattern Matching Algorithms}.

    \item \textbf{Technical Learning \& Continuous Upskilling}
    
    Completed multiple certifications across \textbf{Cybersecurity, Programming, Data Structures, Embedded Systems, IoT, Object-Oriented Design, and Emerging Technologies} through structured online learning.

    \item \textbf{Practical Engineering Development}
    
    Built hands-on systems involving \textbf{biometric authentication, smart monitoring, IoT automation, operating systems simulation, secure workflows, and intelligent software systems}.

\end{itemize}



% ---------- Future Focus Areas ----------

\section{Future Focus Areas}

\textbf{Near-Term Technical Focus}
\begin{itemize}

    \item Strengthening practical expertise in \textbf{Cybersecurity Engineering}, including authentication systems, secure architectures, identity verification, and security-oriented software development.

    \item Expanding capabilities in \textbf{Internet of Things (IoT)} and \textbf{Embedded Systems} through advanced sensor integration, intelligent automation systems, telemetry workflows, and secure device communication.

    \item Deepening software engineering knowledge involving \textbf{Full-Stack Development, scalable application workflows, database systems, APIs, backend engineering, and production-oriented system design}.

    \item Improving problem-solving and computational capabilities through advanced work in \textbf{Data Structures \& Algorithms, Operating Systems, Systems Simulation, and Performance Optimization}.

\end{itemize}

\vspace{4pt}

\textbf{Long-Term Engineering Vision}
\begin{itemize}

    \item Building \textbf{secure, scalable, and intelligent systems} combining cybersecurity, IoT, software engineering, and automation-driven technologies.

    \item Contributing toward \textbf{real-world intelligent infrastructure, secure connected systems, embedded automation, and next-generation engineering solutions}.

    \item Exploring opportunities in \textbf{product engineering, systems architecture, intelligent automation, secure platform development, and innovation-driven technology ecosystems}.

    \item Long-term aspiration to contribute as a \textbf{Software Engineer and Technology Builder} focused on impactful engineering systems and scalable technology products.

\end{itemize}

\vspace{4pt}

\textbf{Continuous Learning Objectives}
\begin{itemize}

    \item Advancing practical skills in \textbf{Cybersecurity, IoT, Full-Stack Engineering, Secure System Design, Cloud Fundamentals, and Intelligent System Development}.

    \item Continuously learning through \textbf{project-based experimentation, engineering implementation, technical research, GitHub development workflows, and structured online education}.

\end{itemize}








% ---------- End Document ----------
\end{document}